\documentclass[11pt,a4paper]{article}

\usepackage[utf8]{inputenc}
\usepackage[T1]{fontenc}
\usepackage{lmodern}
\usepackage{microtype}
\usepackage{geometry}
\geometry{margin=1in}
\usepackage[colorlinks=true,linkcolor=blue,citecolor=blue,urlcolor=blue]{hyperref}
\usepackage{booktabs}
\usepackage{array}
\usepackage{listings}
\usepackage{xcolor}
\usepackage{amsmath}
\usepackage{parskip}
\emergencystretch=1.5em

\lstdefinestyle{code}{
  basicstyle=\ttfamily\small,
  breaklines=true,
  frame=single,
  backgroundcolor=\color{gray!8},
  rulecolor=\color{gray!40},
  xleftmargin=8pt,
  xrightmargin=8pt,
  aboveskip=6pt,
  belowskip=6pt,
}

\title{\textbf{SYNC: Consistent, Cacheable Catch-Up\\of Many HTTP Resources}}
\author{Meet Chauhan\\
\small\href{mailto:meetsc04@gmail.com}{meetsc04@gmail.com}}
\date{October 2026}

\begin{document}

\maketitle

\begin{abstract}
Applications that hold local copies of several server resources must bring all of them up to date, for example when they reconnect. Doing so with one request per resource has three costs that grow with the number of resources: per-request overhead; \emph{torn reads}, in which separate requests observe the server at different moments and combine related resources in a way that never existed; and load concentration, when many clients reconnect at once and the origin computes the same catch-up for each. We describe SYNC, a query format for the HTTP \texttt{QUERY} method (RFC~10008) that addresses all three. A client states the version it holds of each of many resources; the server returns an independent result per resource: a patch, an indication that nothing changed, or the full representation. Resources may have any media type, and versions may be sets of identifiers, compatible with Braid's versioning model. A client can ask for every result to come from one snapshot of the server. It can also let the server redirect it, using the redirection that \texttt{QUERY} defines, to an immutable resource holding the results, which every client in the same state receives, so that ordinary shared caches serve a reconnecting population from one response. We present the design, a reference implementation for Node.js and browsers with 165 tests, a security analysis, and three measurements against the \texttt{braid-http} library, the Mercure hub and plain \texttt{GET} over real sockets. With realistic network and store timing, 82\% of reads made with three parallel requests, with plain \texttt{GET} or with Braid alike, returned combinations of resources that never existed on the server, against none of 300 with SYNC's consistent snapshots, at the same latency. When 100 clients in the same state reconnected through nginx acting as a shared cache, SYNC's shared results reduced the bytes the origin sent from 6.2\,MiB (SYNC without them; 6.6\,MiB from the Mercure hub) to 116\,KiB, against 66\,KiB for per-resource Braid requests that we made cacheable, while each client made 2 requests instead of 50.
\end{abstract}

\section{Introduction}

Applications that keep local state synchronized with a server (mobile applications resuming after being offline, dashboards, editors holding several documents, configuration agents) share a step: the client holds copies of several resources and must bring them current. Most of what it would receive is data it already has.

HTTP offers two standard tools for this. A plain \texttt{GET} returns the full representation every time. A conditional \texttt{GET} returns the full representation or \texttt{304 Not Modified}~\cite{rfc9110}, so a resource that changed in one field is sent again in full. Braid-HTTP~\cite{braid} lets a \texttt{GET} state the version the client holds and returns only the updates since then. All of these work one resource at a time, and catching up many resources with separate requests has three costs:

\begin{enumerate}
  \item \textbf{Overhead.} Each request carries its own headers and often its own connection.
  \item \textbf{Inconsistency.} Separate requests observe the server at different moments. When resources are related (a post and its author, an order and its lines, a configuration and its schema) the client can assemble a combination that never existed. We call this a \emph{torn read}. It occurs even when every individual response is correct.
  \item \textbf{Load concentration.} After an outage or a deployment, many clients reconnect at once, typically from the same state. Unless shared caches can answer them, the origin computes and sends the same catch-up to each.
\end{enumerate}

SYNC is a query format for the \texttt{QUERY} method~\cite{rfc10008} in which a client names many resources, with the version it holds of each, in one request. The server answers each resource independently. On top of this the format provides consistent snapshots, which remove torn reads, and two ways to let ordinary shared caches carry the catch-up: \emph{shared results}, a redirect to an immutable resource holding the results that every client in the same state receives, and \emph{links} to immutable per-resource updates.

An earlier version of this proposal, posted to the IETF HTTP working group, defined a new HTTP method. Discussion there asked how it differed from \texttt{QUERY}; it did not, in any way that matters, and SYNC is now a format for \texttt{QUERY}. Feedback from the Braid authors led to two further changes: versions are sets of identifiers, as in Braid, and representations are no longer restricted to JSON.

\paragraph{Contributions.}
\begin{itemize}
  \item A query format for \texttt{QUERY} with per-resource results for representations of any media type, versions that are sets of identifiers, a JSON and a multipart result format, and a splice patch format for text and binary content (Section~\ref{sec:design}).
  \item \emph{Consistent snapshots} across resources, with an explicit confirmation from the server (Section~\ref{sec:consistency}).
  \item \emph{Shared results}: a use of \texttt{QUERY}'s \texttt{303 (See Other)} redirection with a URI that names the state as well as the query, so that today's shared caches can serve many clients' catch-up; and \emph{links} to immutable updates (Section~\ref{sec:sharing}).
  \item A reference implementation, including stateless, opaque, authenticated URIs for links and shared results, reuse of computed updates across requests, and a client for browsers and Node.js (Section~\ref{sec:impl}).
  \item Three measurements on real sockets against the \texttt{braid-http} library, the Mercure hub and plain \texttt{GET}: single-client overhead, torn reads, and reconnect storms through nginx acting as a shared cache, including the conditions in which SYNC does not help (Section~\ref{sec:eval}).
  \item A security analysis that states which mitigations the implementation enforces (Section~\ref{sec:security}).
\end{itemize}

\section{Related Work}
\label{sec:related}

\paragraph{Conditional GET and delta encoding.}
Entity tags~\cite{rfc9110} give a binary answer. RFC~3229~\cite{rfc3229} adds delta encoding to \texttt{GET}, but the server selects the baseline from its own history; deployment was minimal.

\paragraph{Braid-HTTP.}
Braid~\cite{braid,braidversions} generalizes HTTP towards state synchronization: versions are sets of identifiers forming a directed acyclic graph (\texttt{Version} and \texttt{Parents} headers), updates are patches or snapshots, clients can subscribe, and merge types support concurrent writers. A \texttt{GET} with \texttt{Parents} returns the updates since that version. This is the closest prior art for a single resource, and Braid covers far more than SYNC: live updates and concurrent writers are outside SYNC's scope. Braid requests are per resource; its multiplexing extension~\cite{braidmux} carries many subscriptions over one connection, while each resource is still requested individually. SYNC adopts Braid's model of versions and its \texttt{Version} and \texttt{Parents} fields in multipart results, and is intended to serve as the multi-resource catch-up step for Braid resources.

\paragraph{Mercure and Events Query.}
Mercure~\cite{mercure} is a publish/subscribe hub over Server-Sent Events that resumes from a hub-wide event identifier by replaying every event since then. Events Query~\cite{eventsquery} uses \texttt{QUERY} to return a representation and notifications for one resource, and lists multi-resource delivery and resumption as out of scope.

\paragraph{JMAP, WebDAV synchronization, GraphQL.}
JMAP~\cite{rfc8620} returns changes since a state string within its own object model. WebDAV collection synchronization~\cite{rfc6578} provides a token for one collection. GraphQL batches many reads in one request, but defines neither versions nor updates, and its responses are not consistent across fields unless the server makes them so. SYNC is a generic format for resources identified by URI.

\paragraph{Snapshots and caching.}
Databases avoid torn reads with snapshot isolation or reads ``as of'' a timestamp~\cite{berenson1995}; HTTP has no cross-resource equivalent, and neither has Braid. Content delivery networks absorb load with shared caches~\cite{rfc9111}, which today store responses to \texttt{GET} but not to \texttt{QUERY}.

\section{Design}
\label{sec:design}

\subsection{Why QUERY, Not a New Method}

RFC~9110 asks that a new method be generic and define its caching, conditional and range behavior~\cite{rfc9110}. A safe, idempotent request whose content describes what to return is what \texttt{QUERY} provides, with defined caching (the request content is part of the cache key), conditional requests, and ways to give results a URI~\cite{rfc10008}. A new method would add none of these and would be rejected by parsers and intermediaries that do not know it: Node.js~26 rejects an unknown method in its HTTP parser before application code runs, while it accepts \texttt{QUERY}. SYNC is therefore a media type, \texttt{application/sync-baseline+json}, sent with \texttt{QUERY} to a \emph{sync resource}; a \texttt{POST} fallback serves paths that do not yet allow \texttt{QUERY}.

\subsection{Versions}

A version is a non-empty set of version identifiers, written as a string (one identifier) or an array whose order carries no meaning; versions are equal when they are equal as sets. A version after a single change has one identifier; a version after a merge of concurrent changes can have several, as in Braid. Identifiers are printable ASCII so that they can also travel in Structured Fields~\cite{rfc9651}. A version must identify exactly one state of its resource: the representation associated with a version never changes. This one requirement underlies everything that follows: because the update between two given versions is fixed, a server can reuse it across clients and give it a URI that caches may keep.

\subsection{Request and Per-Resource Results}

\begin{lstlisting}[style=code]
QUERY /sync HTTP/1.1
Content-Type: application/sync-baseline+json
Accept: application/sync-result+json

{ "baselines": { "/users": "a4f2",
                 "/doc.md": ["alice-17", "bob-9"],
                 "/logo.png": null },
  "accept": ["application/merge-patch+json",
             "application/sync-splice+json"],
  "consistent": true, "redirect": true }
\end{lstlisting}

A \emph{baseline} is the version the client holds, or \texttt{null}. Resource names are absolute-path references on the origin of the sync resource, so a request cannot name another origin. Each named resource is subject to the same access control as a \texttt{GET} of it by the same requester, and a resource the requester may not read is reported exactly like one that does not exist.

\begin{lstlisting}[style=code]
{ "results": {
  "/users":    { "status": 200, "from": "a4f2", "to": "c93b",
                 "format": "application/merge-patch+json",
                 "data": { "1": { "email": "new@example.com" } } },
  "/doc.md":   { "status": 200, "from": ["alice-17", "bob-9"],
                 "to": "alice-18",
                 "format": "application/sync-splice+json",
                 "data": { "unit": "codepoint",
                           "splices": [[120, 4, "SYNC"]] } },
  "/logo.png": { "status": 200, "from": null, "to": "l7",
                 "type": "image/png", "encoding": "base64",
                 "data": "iVBORw0KGgo..." } } }
\end{lstlisting}

The response is \texttt{204} when every resource is unchanged and otherwise carries a result per resource: \texttt{200} (an update follows), \texttt{304} (unchanged), \texttt{404} (absent, out of scope, or not readable) or \texttt{409} (baseline not retained, only when the client disables recovery; by default it receives the full representation instead, in the same round trip). A failure for one resource never affects another. This is what makes batching usable: with a single status for the request, one stale version among a hundred would force a retry of everything. Errors in the request itself follow RFC~10008 and are reported as problem details~\cite{rfc9457}.

\subsection{Updates for Any Media Type}

How a representation is carried depends only on its media type: JSON types are JSON values; text types (\texttt{text/*}, XML, JavaScript, and any type with a \texttt{charset}) are UTF-8 text; all others are octets. Updates are JSON Patch~\cite{rfc6902} or JSON Merge Patch~\cite{rfc7396} for JSON types, and a \emph{splice} document for the others: an ordered list of non-overlapping replacements \texttt{[start, deleteCount, insert]}, counted in Unicode code points for text and in octets (with base64 insertions) for binary content. The client lists the patch formats it accepts in preference order; for each resource the server uses the first that applies and can express the change, and sends the full representation when that is not smaller or the media type changed. Recipients validate a splice document completely before applying any of it.

Two result formats are negotiated with \texttt{Accept}. One is the JSON format above. The other is a multipart format~\cite{rfc2046} with one body part per resource, which carries the patch or the representation as octets, so any representation travels without re-encoding; each part has \texttt{Version} and \texttt{Parents} fields in Braid's syntax.

\subsection{Consistent Snapshots}
\label{sec:consistency}

With \texttt{"consistent": true} the server either computes every result from the states of the requested resources at one instant, as if all were read atomically, and says so with \texttt{Sync-Consistent: ?1}, or says \texttt{?0}. It never claims a snapshot it did not take, and a client that requires consistency treats anything but \texttt{?1} as failure. Without this request, results can reflect different moments even within one request, because a server typically reads resources separately; separate requests for the resources have the same problem, more severely (Section~\ref{sec:torn}). The guarantee concerns the states from which results are computed; links and shared results do not weaken it, because each names updates between fixed versions.

\subsection{Shared Results and Links}
\label{sec:sharing}

RFC~10008 lets a server answer a \texttt{QUERY} with \texttt{303 (See Other)}, indicating that the query can be accomplished with a \texttt{GET} of the URI in \texttt{Location}~\cite{rfc10008}. SYNC uses this so that today's shared caches, which store \texttt{GET} responses but not \texttt{QUERY} responses, can carry catch-up. When a request has \texttt{"redirect": true}, the server may respond \texttt{303} with the URI of a \emph{shared result}: a resource whose \texttt{GET} returns exactly the response the request would have received. The URI must identify both the request and the version each result leads to. Two requests therefore receive the same URI only if they would receive the same results, a change to any requested resource yields a different URI, and the content behind a URI never changes, so it can be cached with a long lifetime. This differs from the URI that \texttt{Location} carries in a \texttt{2xx} response to \texttt{QUERY}, which identifies the query and whose results change over time.

The effect is that clients in the same state, which is the common case when a population reconnects after an outage, are sent to one URI, and a shared cache answers all of them from one response. The origin still receives each client's \texttt{QUERY}, but answers it by reading only the current versions of the resources, without computing any update. A server redirects only clients that ask (other clients might not follow a \texttt{303}; Fetch-based clients follow it automatically), never when everything is unchanged, never with a URI longer than intermediaries accept (RFC~9110 recommends at least 8000 octets), and should do so only for results that are the same for every requester.

\emph{Links} apply the same idea per resource: with \texttt{"links": true} the server may replace large content with the URI of an immutable update, which the client retrieves with \texttt{GET}. Shared results suit clients in identical states; links let results for different states share the updates they have in common. Both are subject to per-\texttt{GET} access control, both answer \texttt{404} or \texttt{410} once the versions involved are no longer retained (the client then repeats its request without them), and both must be protected against forgery and should not reveal resource names or versions.

\subsection{Interaction with HTTP}

SYNC requests are \texttt{QUERY} requests: safe, idempotent, and cacheable with the content in the cache key. Results usually depend on the requester, so a server must not let shared caches store them unless they are the same for everyone; the reference server sends \texttt{no-store} by default. A server must treat semantically equivalent request content identically, because caches may normalize it before computing a key. Partial results are signalled with the \texttt{Sync-Delta-Complete} Structured Field. Before applying a patch, a client checks that its baseline equals the version it holds, and if applying any part fails, it leaves its copy unchanged.

\section{Reference Implementation}
\label{sec:impl}

\paragraph{Server.}
A request handler for Node.js, usable as Express middleware or around a plain handler, serves SYNC requests sent with \texttt{QUERY} or \texttt{POST}, links and shared results, and passes every other request on. Its data source is a store with two functions, ``current state'' and ``state at version'', and optionally a snapshot function; each call receives the request context, so the store can authorize per resource. An in-memory store with atomic multi-resource commits and snapshots is included.

Resolution has two steps. \emph{Planning} reads each resource's current version and decides its status and the two versions its update connects. \emph{Materializing} computes the update. A direct response does both. A redirect plans only, and encodes the plan (the request's options and, per resource, its name, status and versions) in the shared result's URI; a \texttt{GET} of that URI re-reads every version involved through the store, with the requester's context, and materializes. This makes the server stateless: any instance behind a load balancer can serve any URI. If any version is no longer retained, or any resource with status \texttt{200} or \texttt{304} is unreadable for the requester, the \texttt{GET} is \texttt{404}.

URIs for links and shared results are deterministic authenticated encryptions of their content: a synthetic IV computed as an HMAC-SHA256 of the plaintext, AES-256 in counter mode, keys derived with HKDF from a configured secret, and constant-time verification. Determinism makes equal requests receive equal URIs, which caching requires; authentication defeats alteration, including bit-flipping of the malleable ciphertext, which the test suite performs. Shared-result payloads are compressed before encryption to keep URIs short (a request naming 100 resources fits in fewer than 3000 characters), and are decompressed only after authentication, with bounded output.

Updates are memoized per store: because a version names one immutable state, the update between two versions is computed once and reused for every client and every link that needs it. Responses that depend on \texttt{Accept-Encoding} carry \texttt{Vary}; the gzip-coded form of a link or shared result has its own strong entity tag; multipart boundaries are derived from the content, so the same results always encode to the same octets.

\paragraph{Client.}
The client is built on \texttt{fetch} and runs in browsers and Node.js. Given resource names it returns their current values (JSON values, strings for text types, byte arrays for others), tracking versions, applying patches, refusing results that do not start from what it holds, following redirects and links, and asking again directly when a shared result or a link has expired or a result does not apply. With \texttt{consistent} it fails unless the server confirms a snapshot. It serializes its state, so a reloaded application resumes with patches. We verified it in Chromium as well as Node.js. Both parts are published as an npm package with TypeScript definitions.

\paragraph{Tests.}
165 automated tests cover, among others, per-resource failure isolation, transports and fallback, RFC~10008 error behavior, problem details, resource-name validation, per-resource authorization, version sets, every media-type class and both result formats, splice validation, consistent snapshots under concurrent writes, links and shared results (determinism, immutability, eviction, authorization of every entry, forgery and bit-flipping, URI length limits, client fallback), content coding and entity tags, and update reuse. We checked the suite by mutation: removing the HMAC check, ignoring the versions in a shared-result URI, skipping re-authorization of unchanged entries, or skipping the boundary check each makes tests fail.

\section{Evaluation}
\label{sec:eval}

All measurements run on real sockets through TCP proxies that count bytes and add a 40\,ms round-trip time, and every client's final state is compared with the server's; any mismatch aborts the run. Resources are JSON objects of 100 items (about 20\,KiB), of which a seeded fraction changes per round. Sizes are in KiB (1024 bytes) and MiB. The real third-party software is the \texttt{braid-http} library~1.5.1 (\texttt{braidify} on the server and its \texttt{fetch} on the client, with Braid range patches using JSON Pointer ranges) and the Mercure.rocks hub~1.1.0 in Docker. Scripts and complete results are in the repository.

\subsection{Single-Client Overhead}

One client brings $N$ resources current after $L$ rounds of change. Table~\ref{tab:bytes} shows bytes as a percentage of full \texttt{GET}, in a \emph{lab} profile (minimal headers, no compression) and a \emph{realistic} profile (bearer token, cookie, browser-style headers, gzip).

\begin{table}[h]
\centering
\footnotesize
\setlength{\tabcolsep}{4pt}
\begin{tabular}{llrrrrrr}
\toprule
 & & \multicolumn{3}{c}{\textbf{Real software}} & \textbf{Model} & & \\
\cmidrule(lr){3-5}
\textbf{Profile} & $(N, L)$ & \textbf{Braid} & \textbf{Braid mux} & \textbf{Mercure} & \textbf{Braid H2} & \textbf{SYNC} & \textbf{GET, KiB} \\
\midrule
Lab       & (1, 1)    & 6.5  & 16.2  & 8.9   & 6.1  & 7.4  & 20.9 \\
Lab       & (10, 1)   & 3.4  & 7.5   & 2.2   & 1.7  & 1.9  & 209.4 \\
Lab       & (100, 1)  & 2.8  & 6.0   & 1.1   & 1.0  & 1.0  & 2094.5 \\
Lab       & (100, 25) & 22.4 & 38.3  & 29.5  & 24.4 & 24.6 & 2096.5 \\
Realistic & (1, 1)    & 34.7 & 73.0  & 50.5  & 31.3 & 27.0 & 5.3 \\
Realistic & (10, 1)   & 22.8 & 38.7  & 10.2  & 8.8  & 3.7  & 53.3 \\
Realistic & (100, 1)  & 20.4 & 33.0  & 4.7   & 6.8  & 1.0  & 532.9 \\
Realistic & (100, 25) & 96.2 & 157.9 & 115.0 & 25.8 & 14.0 & 539.5 \\
\midrule
\multicolumn{2}{l}{Connections at $N=100$} & 100 & 101 & 1 & 1 & 1 & 6 \\
\bottomrule
\end{tabular}
\caption{Total wire bytes (both directions) as a percentage of full \texttt{GET}; lower is better. ``Braid H2'' is a model of a Braid exchange over HTTP/2, which the Node client cannot use over cleartext.}
\label{tab:bytes}
\end{table}

With minimal headers, SYNC and the Mercure hub are close (21.9 and 23.6\,KiB at $N=100$, $L=1$), and real Braid used 58.9\,KiB over 100 connections. With realistic headers SYNC is the smallest in every configuration with more than one resource, but part of that gap is that neither Braid nor Mercure compresses responses by default; the compressing HTTP/2 Braid model isolates the protocol effect. For a single resource SYNC has no consistent advantage, and when many changes accumulate Braid's item-level range patches are smaller than JSON Patch. With a 40\,ms round trip, pooled per-resource \texttt{GET}s of 100 resources took about 800\,ms, against about 100\,ms for SYNC, Mercure and real Braid, the last by opening one connection per resource.

\subsection{Torn Reads}
\label{sec:torn}

A writer commits one transaction every 5\,ms. Each adds a user, adds a post by that user, updates a counter, and removes the oldest user with their post, atomically. In every committed state every post's author is a current user and the counter equals the newest user and post. Clients read \texttt{/users}, \texttt{/posts} and \texttt{/counts} 300 times per approach and check both invariants; every approach reuses its connections. In the \emph{idealized} scenario all delays are fixed, which is the most favourable case for parallel requests because they reach the server at the same instant. In the \emph{realistic} scenario network delays vary by up to 10\,ms and store reads take 1 to 8\,ms (seeded).

\begin{table}[h]
\centering
\footnotesize
\begin{tabular}{lrrr}
\toprule
 & \multicolumn{2}{c}{\textbf{Torn reads (of 300)}} & \\
\cmidrule(lr){2-3}
\textbf{Approach} & \textbf{Idealized} & \textbf{Realistic} & \textbf{Median, realistic} \\
\midrule
\texttt{GET}, three in parallel & 3.0\% & 82.3\% & 58\,ms \\
\texttt{GET}, sequential & 100\% & 100\% & 165\,ms \\
Braid (\texttt{braid-http}), three in parallel & 3.3\% & 82.0\% & 60\,ms \\
SYNC, one request & 0\% & 57.7\% & 57\,ms \\
SYNC, \texttt{consistent} & 0\% & 0\% (95\% CI 0--1.3\%) & 58\,ms \\
\bottomrule
\end{tabular}
\caption{Reads whose combination of resources never existed on the server.}
\label{tab:torn}
\end{table}

Table~\ref{tab:torn} shows the result. Separate requests are torn most of the time once timing varies, and even one SYNC request without a snapshot is torn often, because its server reads the three resources separately. Consistent snapshots removed torn reads entirely at no cost in latency. Separate requests could avoid them only with a mechanism to read ``as of'' a state across resources, which neither HTTP nor Braid defines. The rates depend on the write rate and timing; the comparison shows which approaches can produce states that never existed, not universal percentages.

\subsection{Reconnect Storms}

$K$ clients hold versions of 50 resources and reconnect within one second of each other, five rounds later (41 resources changed). They reach the origin through nginx~1.27 configured as a shared cache with its cache lock, over a 40\,ms round trip, with six connections per client. The origin runs in its own process so that its CPU time can be measured, and every variant starts with a fresh origin and an empty cache. Per-resource \texttt{GET}s and Braid requests are made cacheable (for Braid, \texttt{Cache-Control: public} and \texttt{Vary: Parents}, with the cache keyed on \texttt{Parents}). Nothing is compressed.

\begin{table}[h]
\centering
\footnotesize
\setlength{\tabcolsep}{4pt}
\begin{tabular}{lrrrrr}
\toprule
\textbf{Approach} & \textbf{Origin KiB} & \textbf{Origin req.} & \textbf{Origin CPU} & \textbf{Req./client} & \textbf{Client p50} \\
\midrule
\texttt{GET} (full), cached & 1046 & 50 & 11\,ms & 50 & 399\,ms \\
Braid, cached & 66 & 50 & 18\,ms & 50 & 523\,ms \\
SYNC, inline & 6348 & 100 & 48\,ms & 1 & 94\,ms \\
SYNC, links & 751 & 141 & 56\,ms & 42 & 520\,ms \\
SYNC, shared result & 116 & 101 & 39\,ms & 2 & 183\,ms \\
Mercure hub & 6803 & 100 & n/a & 1 & 100\,ms \\
\bottomrule
\end{tabular}
\caption{100 clients in the same state reconnect through a shared cache. ``Origin'' is what passed the cache. Every SYNC variant computed each of the 41 updates once.}
\label{tab:storm}
\end{table}

Table~\ref{tab:storm} shows the case where all clients were current before the outage. Without a cache, every client receives its own copy: 6.2\,MiB from the SYNC origin and 6.6\,MiB from the Mercure hub. With shared results, the origin answered each \texttt{QUERY} with a small \texttt{303} and the cache served the one shared result to all clients: 116\,KiB at the origin, within a factor of two of cached Braid, while each client made 2 requests instead of 50 and finished in about a third of the time. The \texttt{QUERY} still reaches the origin, so origin traffic grows with $K$, by about half a KiB per client; each \texttt{QUERY} costs reads of current versions and no update computation. With 500 clients (arriving over five seconds), the origin sent 327\,KiB with shared results, against 31\,MiB inline, 33\,MiB from the Mercure hub and 66\,KiB for cached Braid. When clients went offline at different times (five distinct states), shared results sent 259\,KiB from the origin (one shared result per state); with links inside the shared results, which let different states share updates, 180\,KiB; cached Braid, 99\,KiB, at 50 requests per client.

Two observations matter for interpretation. A shared cache absorbs per-resource Braid requests as well once they are made cacheable: caching is not unique to SYNC, and SYNC's advantage here is one request that reveals which resources changed, which a cache cannot provide for separate requests. And the inline variant shows that update reuse alone keeps origin CPU low but not origin bandwidth; only caching removes the bandwidth.

\subsection{Limitations}

The measurements use loopback with emulated latency on one machine: they do not model TLS, packet loss, bandwidth limits or a geographically distributed cache, and the data is synthetic. The third-party software runs in default configurations and is driven by code written for this study; applying Braid's JSON range patches is application code. Braid's Node client cannot use HTTP/2 over cleartext, so HTTP/2 figures come from a model. Shared results help only when clients share a state and results are the same for every requester; for per-user data they must be private and add a round trip without benefit. Only catch-up is measured: Braid and Mercure also provide live push, which SYNC does not.

\section{Security}
\label{sec:security}

The accompanying analysis enumerates attack vectors and states which mitigations the implementation enforces.

\textbf{Authorization and caching.} Authorization applies per named resource, and unreadable resources are indistinguishable from absent ones. Responses, links and shared results that depend on the requester must not be stored by shared caches; the implementation sends \texttt{no-store} or \texttt{private} unless configured otherwise. Every \texttt{GET} of a link or shared result is authorized again for its requester, including entries with status \texttt{304}, which reveal the current version.

\textbf{Links and shared results.} Their URIs can appear in logs, caches and referrers, so they are encrypted and reveal neither resource names nor versions, and authenticated so they cannot be forged or altered. Counter-mode encryption alone is malleable: an attacker who knows the plaintext layout can flip bits that change a version. The authentication tag defeats this. A stateless server replays from the URI only requests it accepted, bounded like the original request. Compressing before encrypting makes the URI length depend on content the requester receives anyway.

\textbf{Replay, probing and malformed patches.} A patch applied to the wrong version corrupts state, so clients check baselines. Any resume mechanism lets a requester learn whether a version is recognized; versions should be unguessable where history is sensitive. Splice documents are validated completely before use, and a failed update leaves the client's copy unchanged.

\textbf{Resource consumption and compression.} One request can name many resources, so limits apply (100 resources, 64\,KiB) and rate limits should count resources. Compressing results that echo request content alongside confidential data invites length-based attacks; \texttt{QUERY} and \texttt{POST} requests of this type require a CORS preflight, which limits cross-origin influence.

The implementation does not provide TLS, authentication, rate limiting or a bound on per-request time.

\section{Discussion and Standardization}
\label{sec:discussion}

\paragraph{Composition.}
Braid, Mercure and Events Query address live synchronization; SYNC is the step a client performs before opening those streams, or instead of them when polling suffices. Its value is greatest when many related resources are involved: one request, results that belong together, and catch-up that caches can share. For a single resource it adds little. SYNC already uses Braid's versions; whether multi-resource catch-up is best specified as a separate format or as an extension of Braid's update model is a question we are pursuing with the Braid authors.

\paragraph{Standardization.}
Because SYNC is a format for an existing method, standardizing it means registering three media types and two header fields, not a method. The draft specification and its open questions (alignment with Braid's patch formats and status codes, a canonical request form for caching, a freshness lifetime for \texttt{303} responses so that caches that learn to store \texttt{QUERY} responses also absorb the requests themselves, and the venue) are in the repository.

\section{Conclusion}

Catching up many resources one request at a time costs overhead, produces combinations of resources that never existed, and concentrates load on the origin when clients reconnect together. SYNC addresses all three with one format for HTTP's \texttt{QUERY} method: per-resource results for any media type, consistent snapshots that removed torn reads in our measurements at no cost in latency, and shared results and links that let ordinary shared caches carry a reconnecting population's catch-up. It reuses \texttt{QUERY}'s semantics, including its redirection, rather than defining new ones, and adopts Braid's versions. Its advantages are conditional and measured, including where it does not help. The implementation, tests, benchmarks, specification and security analysis are available at \url{https://github.com/Meet-1010/sync-http-method}.

\section*{Acknowledgments}

I thank participants of the IETF HTTP working group mailing list, and Michael Toomim of the Braid project, for feedback that changed the design described here.

\begin{thebibliography}{99}\raggedright

\bibitem{rfc9110}
Fielding, R., Nottingham, M., and Reschke, J. ``HTTP Semantics.'' \textit{RFC 9110}, IETF, June 2022.

\bibitem{rfc9111}
Fielding, R., Nottingham, M., and Reschke, J. ``HTTP Caching.'' \textit{RFC 9111}, IETF, June 2022.

\bibitem{rfc10008}
Reschke, J., Snell, J., and Bishop, M. ``The HTTP QUERY Method.'' \textit{RFC 10008}, IETF.

\bibitem{rfc9651}
Nottingham, M. and Kamp, P.-H. ``Structured Field Values for HTTP.'' \textit{RFC 9651}, IETF, September 2024.

\bibitem{rfc9457}
Nottingham, M., Wilde, E., and Dalal, S. ``Problem Details for HTTP APIs.'' \textit{RFC 9457}, IETF, July 2023.

\bibitem{rfc2046}
Freed, N. and Borenstein, N. ``Multipurpose Internet Mail Extensions (MIME) Part Two: Media Types.'' \textit{RFC 2046}, IETF, November 1996.

\bibitem{rfc6902}
Bryan, P. and Nottingham, M. ``JavaScript Object Notation (JSON) Patch.'' \textit{RFC 6902}, IETF, April 2013.

\bibitem{rfc7396}
Hoffman, P. and Snell, J. ``JSON Merge Patch.'' \textit{RFC 7396}, IETF, October 2014.

\bibitem{rfc8620}
Jenkins, N. and Newman, C. ``The JSON Meta Application Protocol (JMAP).'' \textit{RFC 8620}, IETF, July 2019.

\bibitem{rfc6578}
Daboo, C. and Quillaud, A. ``Collection Synchronization for Web Distributed Authoring and Versioning (WebDAV).'' \textit{RFC 6578}, IETF, March 2012.

\bibitem{rfc3229}
Mogul, J., Krishnamurthy, B., Douglis, F., Feldmann, A., Goland, Y., van Hoff, A., and Hellerstein, D. ``Delta encoding in HTTP.'' \textit{RFC 3229}, IETF, January 2002.

\bibitem{berenson1995}
Berenson, H., Bernstein, P., Gray, J., Melton, J., O'Neil, E., and O'Neil, P. ``A Critique of ANSI SQL Isolation Levels.'' In \textit{Proceedings of ACM SIGMOD}, 1995.

\bibitem{braid}
Toomim, M., Little, G., Walker, R., Bellomy, B., and Gentle, J. ``Braid-HTTP: Synchronization for HTTP.'' Internet-Draft \nolinkurl{draft-toomim-httpbis-braid-http-04}, expired.

\bibitem{braidversions}
Toomim, M. ``HTTP Resource Versioning.'' Internet-Draft \nolinkurl{draft-toomim-httpbis-versions-04}, expired.

\bibitem{braidmux}
Braid project. ``Multiplexing.'' \url{https://braid.org/protocol/multiplexing}.

\bibitem{mercure}
Dunglas, K. ``The Mercure Protocol.'' Internet-Draft \nolinkurl{draft-dunglas-mercure-08}.

\bibitem{eventsquery}
Gupta, R. ``HTTP Events Query.'' Internet-Draft \nolinkurl{draft-gupta-httpapi-events-query-03}.

\end{thebibliography}

\end{document}
